%% Response to reviewers, ISASF 2026, Paper ID 132.
%% Red [TO FILL] markers are replaced once the remaining runs finish.
\documentclass[11pt,a4paper]{article}
\usepackage[margin=2.2cm]{geometry}
\usepackage{xcolor}
\usepackage{enumitem}
\usepackage{parskip}
\usepackage[hidelinks]{hyperref}

\newcommand{\rev}[1]{\par\noindent\textbf{Comment.} \textit{#1}\par}
\newcommand{\resp}{\par\noindent\textbf{Response.} }
\newcommand{\change}[1]{\par\noindent\textbf{Changes in the manuscript.} #1\par\medskip}
\newcommand{\tofill}[1]{\textcolor{red}{[TO FILL: #1]}}

\begin{document}

\begin{center}
{\Large\textbf{Response to the Reviewers}}\\[4pt]
ISASF 2026, Paper ID 132\\
\textit{Parameter-Efficient UI Code Generation Using QLoRA on an Extended WebSight Dataset}
\end{center}

We thank the reviewers for their careful reading and constructive comments. We
ran the additional experiments requested where this was possible within the
revision period and revised the manuscript accordingly. Below, each comment is
followed by our response and the corresponding changes. Section, table and
equation numbers refer to the revised manuscript, which remains within the
six-page limit. Two new experiments run through several responses: Qwen2.5-VL-3B
fine-tuned with the same QLoRA recipe, and an evaluation on 100 real web pages
from Design2Code.

\section*{Reviewer 1}

\rev{The dataset extension relies on Qwen2.5-Coder-7B, and a smaller model is
fine-tuned on this synthetic data. This introduces a teacher-student circularity:
the fine-tuned model may learn the biases and artifacts of the larger model, and
its advantage over larger zero-shot models may reflect overfitting to the
teacher's style.}
\resp We agree that code-similarity metrics such as BLEU-4 and CodeBLEU partly
reward agreement with the teacher's style, and we now state this explicitly. Four
results limit the concern. (i) The HTML references are the original WebSight
code, which our teacher did not produce, yet fine-tuning still raises HTML
BLEU-4 from 64.69 to 83.23. (ii) The visual metrics (SSIM, PSNR, LPIPS) compare
renders with the original screenshot, not with the teacher's code, so they do
not reward the teacher's style; we now treat them as the primary measure of
output quality. On them, the fine-tuned model reaches 99\% of the SSIM obtained by
the reference code itself (0.716 vs.\ 0.722) and covers two-thirds of the LPIPS
distance between the zero-shot 3B model and the reference code. (iii) A new
render-based check (below, Reviewer 3, comment 1) shows that the teacher's
targets reproduce their source pages. (iv) On 100 real Design2Code pages, for
which no teacher code exists, fine-tuning still improves every visual metric
over the base model (SSIM $+0.210$, 95\% CI $[0.183, 0.239]$; see comment 2).
\change{New ``Teacher circularity'' paragraph (Sec.~VI); visual metrics declared
primary, with the reference-code ceiling (Sec.~IV-C, Table~I); new render-based
check (Sec.~III-C); Design2Code evaluation (Sec.~V-B); limitation added
(Sec.~VII).}

\rev{The evaluation is performed only on held-out synthetic WebSight pages, which
do not capture the complexity, dynamic states, state management, accessibility
requirements and edge cases of real-world UIs; claiming readiness for
production-grade workflows without real websites (e.g.\ Design2Code) weakens the
contribution.}
\resp We agree and have removed all production-readiness claims, and the
limitations name dynamic states, state management, accessibility and structural
complexity explicitly. We also evaluated the models on real pages from
Design2Code, as suggested. Because the models were trained on single-screen
pages, each page is reduced to its first screen (a 16:9 crop of the top of its
screenshot), which serves as both input and comparison target; we used a random
100 of the 484 pages with the same prompts, decoding and rendering as for
WebSight. All scores drop on real pages, but fine-tuning still helps. The
fine-tuned 2B model improves on its base model by 0.210 SSIM $[0.183, 0.239]$,
3.2\,dB PSNR and 0.122 LPIPS (paired bootstrap over pages, all $p<0.001$),
renders 100/93/99\% of its HTML/JSX/Vue outputs (base model: 96/13/90\%), and is
statistically tied with the fine-tuned 3B model on all three metrics
($p \ge 0.08$). Its HTML reaches SSIM 0.716 on the 94 pages whose real HTML
renders, against 0.730 for the real HTML itself. The zero-shot 3B model, about
twice as slow, was run on the first 25 of these pages. There, the fine-tuned 2B
model is statistically tied with it (SSIM $-0.021$ $[-0.071, 0.032]$, LPIPS
$-0.012$ $[-0.052, 0.027]$), so the LPIPS advantage we observed on WebSight is
not confirmed on real pages; with 25 pages this comparison has little
statistical power, and we report it as such.
The evaluation also exposed a limitation that we now report: the fine-tuned
model writes about 600 tokens per page, the scale of its WebSight training
targets, whereas the real HTML has a median of 14,805 tokens, so complex pages
are simplified.
\change{Abstract; Sec.~IV-C (protocol); new Sec.~V-B with Table~III; Sec.~VI;
Sec.~VII (limitations and future work).}

\rev{Comparing a fine-tuned 2B model with zero-shot baselines is unbalanced. A
proper baseline is the 3B model (or a similar model) fine-tuned on the same data
with LoRA/QLoRA; otherwise the gains cannot be attributed to QLoRA rather than to
in-domain supervised fine-tuning.}
\resp We agree with this distinction and made two changes. First, we no longer
attribute the gains to QLoRA itself: the revised text states that the results show
the benefit of in-domain fine-tuning, which QLoRA makes affordable on an 8\,GB
GPU, and not that QLoRA is better than other adaptation methods. Second, we
fine-tuned Qwen2.5-VL-3B on exactly the same data with the same QLoRA recipe
(rank 64, $\alpha=128$, learning rate $10^{-4}$, same data and seed). To fit the
larger model on the same 8\,GB GPU we used bf16 precision, a paged 8-bit AdamW
optimizer and 16-bit frozen embeddings, and trained for one epoch (3,152 steps,
about 34 hours; three epochs would have taken more than four days). Its training loss was lower than
that of the 2B model in all 32 windows of 100 steps (by 6--14\%, median 10\%),
reaching 0.053 vs.\ 0.061 over the last 100 steps of the epoch.\\
On the test set, fine-tuning closes most of the 3B model's gap to the fine-tuned
2B model (BLEU-4 from 54.31 to 78.41, LPIPS from 0.563 to 0.524) but does not
overtake it. The two are tied on SSIM ($-0.0004$, 95\% CI $[-0.005, 0.004]$,
$p=0.87$), and the 2B model is slightly better on LPIPS ($-0.009$
$[-0.015, -0.003]$), PSNR ($+0.15$\,dB $[0.01, 0.30]$), BLEU-4 ($+2.3$
$[1.6, 3.0]$) and CodeBLEU ($+0.022$ $[0.015, 0.028]$). On Design2Code the two
fine-tuned models are tied on all visual metrics. The fine-tuned 3B model needs
51\% more GPU memory and 15\% more time per sample (3.30 vs.\ 2.18\,GiB; 12.9 vs.\
11.2\,s at batch size 8). Training for one epoch instead of three favours the 2B
model in training budget, which we state; at this budget, the larger model brings
no gain.
\change{Sec.~IV-A (new model configuration); Tables~I and II (new QLoRA (3B)
column); Sec.~V-A (new comparison paragraph and summary); Sec.~V-B; Sec.~VI
(new ``Model size'' paragraph); Sec.~VII (single training run, one epoch, BF16).}

\section*{Reviewer 3}

We thank the reviewer for verifying the internal consistency of our headline
figures and for the detailed suggestions.

\rev{(1) The reference code is circular; BLEU and CodeBLEU reward imitating the
teacher. Please treat the visual metrics as primary, where the 3B model is already
close (SSIM 0.711 vs.\ 0.716), and verify the dataset by rendering each generated
JSX/Vue file and comparing it with its source HTML screenshot.}
\resp On circularity, please see our response to Reviewer~1, comment~1. The visual
metrics are now primary (listed first in Table~I), and we report paired bootstrap
tests for every comparison (comment~4). The reviewer is right that SSIM does not
separate the two models: the difference is $+0.005$ (95\% CI $[-0.004, 0.014]$,
$p=0.33$), and on Vue the zero-shot 3B model is significantly better. SSIM is
close to its ceiling for both models (0.722 for the reference code itself). The
fine-tuned model is, however, significantly better on LPIPS ($-0.048$, CI
$[-0.057, -0.040]$) and PSNR ($+0.53$\,dB, CI $[0.29, 0.77]$). We have rewritten
the claims accordingly.\\
For the dataset check, we rendered the original HTML, the JSX target and the Vue
target of all 496 test images and of a random sample of 1,504 training and
validation images with the evaluation harness, and compared each target's render
with the render of its source HTML. JSX and Vue targets rendered in 99.5\% and
100\% of cases, with a median SSIM to the source render of 0.90 and 0.96. Lower
scores came mainly from page-level styling (e.g.\ a dark \texttt{<body>}
background that a mounted component covers only up to its content height) and
from randomly sampled placeholder images, not from missing content: low- and
high-scoring targets had the same ratio of elements to their source page.
\change{Sec.~III-C (render-based check), Sec.~IV-C, Table~I, Sec.~V-A.}

\rev{(2) The baselines are unequal. Please add Qwen2.5-VL-3B with the same QLoRA
training, few-shot prompting baselines, and an existing WebSight-trained model.}
\resp The QLoRA-trained Qwen2.5-VL-3B is now included; please see our response to
Reviewer~1, comment~3. We were not able to add the other two baselines within the
revision period on our hardware, and we explain why. A few-shot prompt for this
task must contain at least one example screenshot and its full code (about 1,500
tokens per example) in addition to the query image. The zero-shot 3B model in
16-bit precision already needs more memory than our 8\,GB GPU provides, so a
longer few-shot context was not practical within the revision period, and the
results would also depend on the choice of examples. The WebSight-trained models
we are aware of (e.g.\ the model released with WebSight and Design2Code-18B) have
7B parameters or more and generate HTML only, so they cannot be run on our
hardware and cannot produce the JSX and Vue outputs evaluated here. We list both
comparisons as future work.
\change{Sec.~IV-A, Sec.~V-A, Sec.~VII.}

\rev{(3) The results are from synthetic data only. WebSight pages are
template-like, so a random split may leak near-duplicates. Please de-duplicate,
and evaluate on the real-website Design2Code benchmark.}
\resp The split is made by image, so the three tasks of an image always fall in
the same split. We computed, for every test page, the word 5-gram Jaccard
similarity of its HTML to the closest training page. No test page duplicated a
training page, one exceeded the common near-duplicate threshold of 0.8, and 28
exceeded 0.5. Removing these 28 pages changed SSIM, LPIPS and CodeBLEU by less
than 0.01, PSNR by less than 0.2\,dB and BLEU-4 by less than 0.5 points for every
model, and did not change the ranking of the models. For Design2Code, please see
our response to Reviewer~1, comment~2.
\change{Sec.~III-D; Design2Code: Sec.~IV-C, Sec.~V-B, Table~III.}

\rev{(4) No variance is reported. Please give bootstrap confidence intervals or
results over several seeds. Please also explain why the 4-bit model runs faster
than the FP16 base model (shorter outputs?) and report the number of generated
tokens.}
\resp We now report 95\% confidence intervals and two-sided $p$-values from a
paired bootstrap over the 496 test images (1,000 resamples; 500 for CodeBLEU),
in which the three tasks of an image are resampled together and all models are
compared on the same resamples. Repeated training runs with different seeds were
not feasible on our hardware within the revision period, and we list this as a
limitation.\\
On latency, the reviewer's intuition explains our original numbers, and a
controlled re-run changed our conclusion. Generation is batched, and a batch
finishes only when its longest output does. In the original runs, batch time
correlated with the longest output in the batch ($r=0.86$--$0.91$), the base
model reached the 1,536-token limit in 3.7\% of its outputs, and the models ran
with different batch sizes (20 to 40). We therefore timed all four models on the
same 24 test samples with batch size 8, where no output reached the limit. Per
sample, the fine-tuned 2B model needs 11.2\,s and 2.18\,GiB of peak GPU memory,
the base model 6.4\,s and 4.59\,GiB, the fine-tuned 3B model 12.9\,s and
3.30\,GiB, and the zero-shot 3B model 22.9\,s and 6.47\,GiB, with mean outputs of
433, 389, 407 and 347 tokens. Under equal conditions, the 4-bit model is thus
1.75$\times$ slower than the FP16 base model, which we attribute to the
dequantization of its weights and the separately applied adapter, while it halves
the memory and is twice as fast as the zero-shot 3B model. We removed the claim
that the fine-tuned model is faster than its base model.
\change{Sec.~IV-C, Sec.~V (Efficiency paragraph), Sec.~VII.}

\rev{(5) There are internal inconsistencies: the maximum sequence length is 2048
in Table IV but 1,536 in the Limitations. The base model's 12.1\% JSX render
rate may partly reflect the rendering harness (unresolved imports) rather than
the model itself.}
\resp The two numbers refer to different quantities, which we failed to make
clear: 2,048 tokens is the maximum training sequence length (image, prompt and
code), and 1,536 is the maximum number of generated tokens at inference, which
covers 99.93\% of the test references. Sec.~IV-B now states the first and
Sec.~IV-C the second, and the limitations give the correct truncation rates
(3.7\% of the base model's outputs, 1.1\% of the zero-shot 3B model's and none of
the fine-tuned model's). The hyperparameter table, which repeated the text of
Sec.~IV-B, was removed for space.\\
On the render rate, the reviewer is right. 437 of the base model's 496 JSX outputs
import a UI library such as react-bootstrap, despite the Tailwind-only
instruction (19 for the zero-shot 3B model, none for the fine-tuned model), and
our harness removes import statements. We added a lenient variant that replaces
imported library components with plain \texttt{<div>} stubs: 87.7\% of the base
model's JSX outputs then render (SSIM 0.638). The strict render gap therefore
mostly reflects instruction following, and with the lenient harness the overall
SSIM gap between the base and fine-tuned models narrows to 0.027. We report both
settings and removed the earlier ``+40.7\% visual accuracy'' claim. The same
library imports explain the base model's 13\% JSX render rate on Design2Code.
\change{Sec.~IV-B, Sec.~IV-C, Sec.~V, Sec.~V-A, Sec.~VII.}

\rev{(6) Several citations are wrong. GPT-3, GPT-4, Transformer, ViT, CLIP and
Flamingo are cited as UI-to-code research. Design2Code is by Si et al.\ (NAACL
2025), and UICoder is by Wu et al. Related work also misses Web2Code, WebCode2M
and DCGen.}
\resp Thank you; all corrected. The foundation models are now cited as such, and
UI-to-code research is cited separately (pix2code, REMAUI, WebSight,
Design2Code). Design2Code is now Si et al.\ (NAACL 2025) and UICoder Wu et al.\
(NAACL 2024), and Web2Code, WebCode2M and DCGen are discussed in Related Work,
together with a paragraph on how our work differs from them. While checking the
bibliography we also corrected the Qwen2-VL and PEFT-survey entries, updated
arXiv citations to their published venues (e.g.\ pix2code, EICS 2018), and added
the missing citations for the BLEU, CodeBLEU and LPIPS metrics.
\change{Sec.~I, Sec.~II, references.}

\section*{Questionnaire items}

\textbf{Differentiation from existing literature (both reviewers).} Related Work
now ends with a paragraph contrasting our work with Design2Code, Web2Code,
WebCode2M, DCGen and UICoder: these systems produce HTML/CSS or platform UI code
and rely on large or proprietary models or on training at scale, whereas our work
provides validated React and Vue targets and adapts a small open model on consumer
hardware.

\textbf{Research objectives (Reviewer~1).} The contributions in Sec.~I are now
stated as objectives, each followed by what was done to meet it.

\textbf{Assumptions (both reviewers).} We now justify the image resolution
(256--512 visual tokens per screenshot, Sec.~III-D), the generation limit
(Sec.~IV-C) and the use of teacher-generated targets (Sec.~III-C and Sec.~VI).
While doing so we also corrected our description of the image pipeline: training
used the full-resolution screenshots, downscaled by the model's image processor,
and the 640$\times$360 copies were used for evaluation.

\textbf{Changed numbers.} While re-running the evaluation, we found that the
data-flow component of CodeBLEU depends on Python's hash seed, which moved JSX
CodeBLEU by up to 0.01 between runs. All CodeBLEU values are now computed with a
fixed seed; the overall values change by at most 0.003 (e.g.\ 0.732 to 0.735 for
the fine-tuned model), and no comparison changes. Latency and memory figures are
now those of the controlled re-run (Reviewer~3, comment~4).

\end{document}
